\QHeading{第1题　直接写得数。}
\par\noindent


\CalcBlock[方法一]{1}{\begin{aligned}[t]&\phantom{={}}\color{black}(15+25)\times2\\&\color{answerred}=40\times2\\&\color{answerred}=80\end{aligned}}
\par
\MethodBlock{方法二：乘法分配律}{分别求出两部分的积，再合起来。}{\begin{aligned}[t]&\phantom{={}}(15+25)\times 2\\
&=15\times 2+25\times 2\\
&=30+50\\
&=80\end{aligned}}


\CalcBlock{2}{\begin{aligned}[t]&\phantom{={}}\color{black}(80-20)\div3\\&\color{answerred}=60\div3\\&\color{answerred}=20\end{aligned}}
\CalcBlock[方法一]{3}{\begin{aligned}[t]&\phantom{={}}\color{black}5\times(12+18)\\&\color{answerred}=5\times30\\&\color{answerred}=150\end{aligned}}
\MethodBlock{方法二：乘法分配律}{把括号中的两部分分别相乘，再相加或相减。}{\begin{aligned}[t]&\phantom{={}}5\times (12+18)\\
&=5\times 12+5\times 18\\
&=60+90\\
&=150\end{aligned}}

\CalcBlock{4}{\begin{aligned}[t]&\phantom{={}}\color{black}36\div(4+5)\\&\color{answerred}=36\div9\\&\color{answerred}=4\end{aligned}}
\CalcBlock{5}{\begin{aligned}[t]&\phantom{={}}\color{black}(25+15)\div8\\&\color{answerred}=40\div8\\&\color{answerred}=5\end{aligned}}
\CalcBlock[方法一]{6}{\begin{aligned}[t]&\phantom{={}}\color{black}8\times(30-15)\\&\color{answerred}=8\times15\\&\color{answerred}=120\end{aligned}}
\MethodBlock{方法二：乘法分配律}{把括号中的两部分分别相乘，再相加或相减。}{\begin{aligned}[t]&\phantom{={}}8\times (30-15)\\
&=8\times 30-8\times 15\\
&=240-120\\
&=120\end{aligned}}

\CalcBlock{7}{\begin{aligned}[t]&\phantom{={}}\color{black}60\div(15-9)\\&\color{answerred}=60\div6\\&\color{answerred}=10\end{aligned}}
\CalcBlock{8}{\begin{aligned}[t]&\phantom{={}}\color{black}(36+24)\div5\\&\color{answerred}=60\div5\\&\color{answerred}=12\end{aligned}}
\CalcBlock{9}{\begin{aligned}[t]&\phantom{={}}\color{black}30\times(100\div5)\\&\color{answerred}=30\times20\\&\color{answerred}=600\end{aligned}}
\par\noindent

\QHeading{第2题　解决问题。（先画图整理条件和问题，再解答。）}

\Problem{①芳芳每分钟走80米，欣欣每分钟走75米，两人同时从家出发，5分钟后同时到达少年宫。两家相距多少米？}
\par




\MethodTitle{方法一}

\par\noindent\makebox[\linewidth][c]{
\begin{tikzpicture}[x=1bp,y=-1bp,color=black,draw=black,text=black]
\begin{scope}\clip (303.0,310.67) rectangle (572.33,391.2);
% 高清原题重排：PDF 26，原书 21。不含解答。
\hdhead{21}{解决问题——信息窗3}
\hdsection{1}{145}{直接写得数。}
\q{96}{192}{193}{13}{\((15+25)\times2=\)}
\q{350}{192}{193}{13}{\((80-20)\div3=\)}
\q{597}{192}{193}{13}{\(5\times(12+18)=\)}
\q{96}{236}{193}{13}{\(36\div(4+5)=\)}
\q{350}{236}{193}{13}{\((25+15)\div8=\)}
\q{597}{236}{193}{13}{\(8\times(30-15)=\)}
\q{96}{280}{193}{13}{\(60\div(15-9)=\)}
\q{350}{280}{193}{13}{\((36+24)\div5=\)}
\q{597}{280}{193}{13}{\(30\times(100\div5)=\)}
\hdsection{2}{334}{解决问题。（先画图整理条件和问题，再解答。）}



\q{462}{475}{105}{13.5}{芳芳家}
\q{615}{475}{105}{13.5}{少年宫}
\q{746}{475}{105}{13.5}{欣欣家}
\begin{scope}[scale=.6666667,draw=black,line width=.65pt]
\draw (495,509)--(779,509);
\fill (495,509) circle (1.9) (649,509) circle (1.9) (779,509) circle (1.9);
\end{scope}
\q{95}{597}{31}{13.5}{②}
\q{128}{598}{716}{13.5}{甲、乙两辆汽车同时从两地相对开出，甲车速度为56千米/时，乙车}
\q{128}{640}{716}{13.5}{速度为64千米/时，4小时后相遇。两地相距多少千米？}
\q{95}{783}{31}{13.5}{③}
\q{128}{784}{716}{13.5}{甲、乙两个工程队一起铺路，甲队每天铺140米，乙队每天铺150米，}
\q{128}{826}{716}{13.5}{两队一起干了5天，一共铺路多少米？}
\q{95}{972}{31}{13.5}{④}
\q{128}{973}{716}{13.5}{两艘船分别从A、B两港口相对开出，甲船的速度为32千米/时，乙}
\q{128}{1015}{716}{13.5}{船的速度为28千米/时，行驶5小时后还相距80千米。A、B两港口}
\q{128}{1057}{716}{13.5}{相距多少千米？}

\begin{scope}[color=answerred,draw=answerred,text=answerred]\ans{165.33}{137.33}{62.00}{9}{\linespread{1}\selectfont\setlength{\jot}{-2pt}\(\begin{aligned}[t]&=40\times2\\&=80\end{aligned}\)}
\ans{333.33}{137.33}{62.00}{9}{\linespread{1}\selectfont\setlength{\jot}{-2pt}\(\begin{aligned}[t]&=60\div3\\&=20\end{aligned}\)}
\ans{504.67}{137.33}{62.00}{9}{\linespread{1}\selectfont\setlength{\jot}{-2pt}\(\begin{aligned}[t]&=5\times30\\&=150\end{aligned}\)}
\ans{165.33}{167.33}{62.00}{9}{\linespread{1}\selectfont\setlength{\jot}{-2pt}\(\begin{aligned}[t]&=36\div9\\&=4\end{aligned}\)}
\ans{333.33}{167.33}{62.00}{9}{\linespread{1}\selectfont\setlength{\jot}{-2pt}\(\begin{aligned}[t]&=40\div8\\&=5\end{aligned}\)}
\ans{504.67}{167.33}{62.00}{9}{\linespread{1}\selectfont\setlength{\jot}{-2pt}\(\begin{aligned}[t]&=8\times15\\&=120\end{aligned}\)}
\ans{165.33}{196.00}{62.00}{9}{\linespread{1}\selectfont\setlength{\jot}{-2pt}\(\begin{aligned}[t]&=60\div6\\&=10\end{aligned}\)}
\ans{333.33}{196.00}{62.00}{9}{\linespread{1}\selectfont\setlength{\jot}{-2pt}\(\begin{aligned}[t]&=60\div5\\&=12\end{aligned}\)}
\ans{504.67}{196.00}{62.00}{9}{\linespread{1}\selectfont\setlength{\jot}{-2pt}\(\begin{aligned}[t]&=30\times20\\&=600\end{aligned}\)}

\draw[answerred,line width=.9pt,->] (331.333,339.333)--(430.000,339.333);
\draw[answerred,line width=.9pt,->] (519.333,339.333)--(430.667,339.333);
\ans{350.67}{350.67}{80.67}{10}{400米}
\ans{454.00}{350.67}{73.33}{10}{375米}
\ans{330.00}{368.00}{212.00}{9.5}{芳芳：80米/分；欣欣：75米/分}
\ans{86.00}{451.33}{229.33}{11}{\linespread{1}\selectfont\setlength{\jot}{-2pt}\(\begin{aligned}[t]&\phantom{={}}(56+64)\times4\\&=120\times4\\&=480\text{（千米）}\end{aligned}\)\\答：两地相距480千米。}
\draw[answerred,line width=.8pt,->] (340.667,480.000)--(427.780,480.000);
\draw[answerred,line width=.8pt,->] (527.333,480.000)--(427.780,480.000);
\ans{333.33}{462.67}{37.33}{10}{甲}
\ans{520.00}{462.67}{33.33}{10}{乙}
\ans{414.00}{460.67}{52.67}{10}{相遇}
\ans{350.00}{489.33}{80.00}{9.5}{224千米}
\ans{448.00}{489.33}{80.00}{9.5}{256千米}
\ans{357.33}{504.00}{178.67}{9.5}{两车都行驶4小时}
\ans{86.00}{574.00}{244.00}{11}{\linespread{1}\selectfont\setlength{\jot}{-2pt}\(\begin{aligned}[t]&\phantom{={}}(140+150)\times5\\&=290\times5\\&=1450\text{（米）}\end{aligned}\)\\答：一共铺路1450米。}
\ans{360.67}{567.33}{170.67}{9.5}{甲：每天140米，铺5天}
\ans{360.67}{596.67}{170.67}{9.5}{乙：每天150米，铺5天}
\draw[answerred,line width=.8pt] (361.333,588.000)--(464.000,588.000);
\draw[answerred,line width=.8pt] (361.333,617.333)--(471.333,617.333);
\ans{473.33}{580.00}{74.00}{10}{700米}
\ans{480.00}{609.33}{71.33}{10}{750米}
\ans{86.00}{728.00}{245.33}{10.5}{\linespread{1}\selectfont\setlength{\jot}{-2pt}\(\begin{aligned}[t]&\phantom{={}}(32+28)\times5+80\\&=60\times5+80\\&=300+80\\&=380\text{（千米）}\end{aligned}\)\\答：两港口相距380千米。}
\draw[answerred,line width=.8pt] (343.333,751.333)--(536.667,751.333);
\draw[answerred,line width=.9pt,->] (343.333,751.333)--(424.737,751.333);
\draw[answerred,line width=.9pt,->] (536.667,751.333)--(465.439,751.333);
\ans{338.67}{733.33}{36.67}{10}{\(A\)}
\ans{529.33}{733.33}{30.00}{10}{\(B\)}
\ans{348.00}{760.00}{80.67}{9}{160千米}
\ans{425.33}{760.00}{48.67}{9}{80千米}
\ans{478.00}{760.00}{76.67}{9}{140千米}
\ans{388.00}{774.00}{168.67}{9}{两船各行驶5小时}
\draw[answerred,line width=.6pt] (361.333,585.333)--(361.333,590.667);
\draw[answerred,line width=.6pt] (381.867,585.333)--(381.867,590.667);
\draw[answerred,line width=.6pt] (402.400,585.333)--(402.400,590.667);
\draw[answerred,line width=.6pt] (422.933,585.333)--(422.933,590.667);
\draw[answerred,line width=.6pt] (443.467,585.333)--(443.467,590.667);
\draw[answerred,line width=.6pt] (464.000,585.333)--(464.000,590.667);
\draw[answerred,line width=.6pt] (361.333,614.667)--(361.333,620.000);
\draw[answerred,line width=.6pt] (383.333,614.667)--(383.333,620.000);
\draw[answerred,line width=.6pt] (405.333,614.667)--(405.333,620.000);
\draw[answerred,line width=.6pt] (427.333,614.667)--(427.333,620.000);
\draw[answerred,line width=.6pt] (449.333,614.667)--(449.333,620.000);
\draw[answerred,line width=.6pt] (471.333,614.667)--(471.333,620.000);
\end{scope}
\end{scope}\pgfresetboundingbox\path[use as bounding box] (303.0,310.67) rectangle (572.33,391.2);
\end{tikzpicture}}\par
\Answer{\Formula{\begin{aligned}[t]&\phantom{={}}(80+75)\times5\\&=155\times5\\&=775\text{（米）}\end{aligned}}\par 答：两家相距775米。}

\MethodTitle{方法二：分别计算两部分，再求和}
\Answer{不先合并每份的数量，而是分别求出两部分的总数，再相加。}
\Answer{\Formula{\begin{aligned}[t]&\phantom{={}}80\times 5+75\times 5\\
&=400+375\\
&=775\text{（米）}\end{aligned}}}
\Answer{答：两家相距775米。}
\Problem{②甲、乙两辆汽车同时从两地相对开出，甲车速度为56千米/时，乙车速度为64千米/时，4小时后相遇。两地相距多少千米？}



\MethodTitle{方法一}

\par\noindent\makebox[\linewidth][c]{
\begin{tikzpicture}[x=1bp,y=-1bp,color=black,draw=black,text=black]
\begin{scope}\clip (328.33,454.67) rectangle (558.33,519.0);
% 高清原题重排：PDF 26，原书 21。不含解答。
\hdhead{21}{解决问题——信息窗3}
\hdsection{1}{145}{直接写得数。}
\q{96}{192}{193}{13}{\((15+25)\times2=\)}
\q{350}{192}{193}{13}{\((80-20)\div3=\)}
\q{597}{192}{193}{13}{\(5\times(12+18)=\)}
\q{96}{236}{193}{13}{\(36\div(4+5)=\)}
\q{350}{236}{193}{13}{\((25+15)\div8=\)}
\q{597}{236}{193}{13}{\(8\times(30-15)=\)}
\q{96}{280}{193}{13}{\(60\div(15-9)=\)}
\q{350}{280}{193}{13}{\((36+24)\div5=\)}
\q{597}{280}{193}{13}{\(30\times(100\div5)=\)}
\hdsection{2}{334}{解决问题。（先画图整理条件和问题，再解答。）}
\q{95}{378}{31}{13.5}{①}
\q{128}{379}{716}{13.5}{芳芳每分钟走80米，欣欣每分钟走75米，两人同时从家出发，5分}
\q{128}{421}{716}{13.5}{钟后同时到达少年宫。两家相距多少米？}
\q{462}{475}{105}{13.5}{芳芳家}
\q{615}{475}{105}{13.5}{少年宫}
\q{746}{475}{105}{13.5}{欣欣家}
\begin{scope}[scale=.6666667,draw=black,line width=.65pt]
\draw (495,509)--(779,509);
\fill (495,509) circle (1.9) (649,509) circle (1.9) (779,509) circle (1.9);
\end{scope}



\q{95}{783}{31}{13.5}{③}
\q{128}{784}{716}{13.5}{甲、乙两个工程队一起铺路，甲队每天铺140米，乙队每天铺150米，}
\q{128}{826}{716}{13.5}{两队一起干了5天，一共铺路多少米？}
\q{95}{972}{31}{13.5}{④}
\q{128}{973}{716}{13.5}{两艘船分别从A、B两港口相对开出，甲船的速度为32千米/时，乙}
\q{128}{1015}{716}{13.5}{船的速度为28千米/时，行驶5小时后还相距80千米。A、B两港口}
\q{128}{1057}{716}{13.5}{相距多少千米？}

\begin{scope}[color=answerred,draw=answerred,text=answerred]\ans{165.33}{137.33}{62.00}{9}{\linespread{1}\selectfont\setlength{\jot}{-2pt}\(\begin{aligned}[t]&=40\times2\\&=80\end{aligned}\)}
\ans{333.33}{137.33}{62.00}{9}{\linespread{1}\selectfont\setlength{\jot}{-2pt}\(\begin{aligned}[t]&=60\div3\\&=20\end{aligned}\)}
\ans{504.67}{137.33}{62.00}{9}{\linespread{1}\selectfont\setlength{\jot}{-2pt}\(\begin{aligned}[t]&=5\times30\\&=150\end{aligned}\)}
\ans{165.33}{167.33}{62.00}{9}{\linespread{1}\selectfont\setlength{\jot}{-2pt}\(\begin{aligned}[t]&=36\div9\\&=4\end{aligned}\)}
\ans{333.33}{167.33}{62.00}{9}{\linespread{1}\selectfont\setlength{\jot}{-2pt}\(\begin{aligned}[t]&=40\div8\\&=5\end{aligned}\)}
\ans{504.67}{167.33}{62.00}{9}{\linespread{1}\selectfont\setlength{\jot}{-2pt}\(\begin{aligned}[t]&=8\times15\\&=120\end{aligned}\)}
\ans{165.33}{196.00}{62.00}{9}{\linespread{1}\selectfont\setlength{\jot}{-2pt}\(\begin{aligned}[t]&=60\div6\\&=10\end{aligned}\)}
\ans{333.33}{196.00}{62.00}{9}{\linespread{1}\selectfont\setlength{\jot}{-2pt}\(\begin{aligned}[t]&=60\div5\\&=12\end{aligned}\)}
\ans{504.67}{196.00}{62.00}{9}{\linespread{1}\selectfont\setlength{\jot}{-2pt}\(\begin{aligned}[t]&=30\times20\\&=600\end{aligned}\)}
\ans{86.00}{308.67}{223.33}{11}{\linespread{1}\selectfont\setlength{\jot}{-2pt}\(\begin{aligned}[t]&\phantom{={}}(80+75)\times5\\&=155\times5\\&=775\text{（米）}\end{aligned}\)\\答：两家相距775米。}
\draw[answerred,line width=.9pt,->] (331.333,339.333)--(430.000,339.333);
\draw[answerred,line width=.9pt,->] (519.333,339.333)--(430.667,339.333);
\ans{350.67}{350.67}{80.67}{10}{400米}
\ans{454.00}{350.67}{73.33}{10}{375米}
\ans{330.00}{368.00}{212.00}{9.5}{芳芳：80米/分；欣欣：75米/分}

\draw[answerred,line width=.8pt,->] (340.667,480.000)--(427.780,480.000);
\draw[answerred,line width=.8pt,->] (527.333,480.000)--(427.780,480.000);
\ans{333.33}{462.67}{37.33}{10}{甲}
\ans{520.00}{462.67}{33.33}{10}{乙}
\ans{414.00}{460.67}{52.67}{10}{相遇}
\ans{350.00}{489.33}{80.00}{9.5}{224千米}
\ans{448.00}{489.33}{80.00}{9.5}{256千米}
\ans{357.33}{504.00}{178.67}{9.5}{两车都行驶4小时}
\ans{86.00}{574.00}{244.00}{11}{\linespread{1}\selectfont\setlength{\jot}{-2pt}\(\begin{aligned}[t]&\phantom{={}}(140+150)\times5\\&=290\times5\\&=1450\text{（米）}\end{aligned}\)\\答：一共铺路1450米。}
\ans{360.67}{567.33}{170.67}{9.5}{甲：每天140米，铺5天}
\ans{360.67}{596.67}{170.67}{9.5}{乙：每天150米，铺5天}
\draw[answerred,line width=.8pt] (361.333,588.000)--(464.000,588.000);
\draw[answerred,line width=.8pt] (361.333,617.333)--(471.333,617.333);
\ans{473.33}{580.00}{74.00}{10}{700米}
\ans{480.00}{609.33}{71.33}{10}{750米}
\ans{86.00}{728.00}{245.33}{10.5}{\linespread{1}\selectfont\setlength{\jot}{-2pt}\(\begin{aligned}[t]&\phantom{={}}(32+28)\times5+80\\&=60\times5+80\\&=300+80\\&=380\text{（千米）}\end{aligned}\)\\答：两港口相距380千米。}
\draw[answerred,line width=.8pt] (343.333,751.333)--(536.667,751.333);
\draw[answerred,line width=.9pt,->] (343.333,751.333)--(424.737,751.333);
\draw[answerred,line width=.9pt,->] (536.667,751.333)--(465.439,751.333);
\ans{338.67}{733.33}{36.67}{10}{\(A\)}
\ans{529.33}{733.33}{30.00}{10}{\(B\)}
\ans{348.00}{760.00}{80.67}{9}{160千米}
\ans{425.33}{760.00}{48.67}{9}{80千米}
\ans{478.00}{760.00}{76.67}{9}{140千米}
\ans{388.00}{774.00}{168.67}{9}{两船各行驶5小时}
\draw[answerred,line width=.6pt] (361.333,585.333)--(361.333,590.667);
\draw[answerred,line width=.6pt] (381.867,585.333)--(381.867,590.667);
\draw[answerred,line width=.6pt] (402.400,585.333)--(402.400,590.667);
\draw[answerred,line width=.6pt] (422.933,585.333)--(422.933,590.667);
\draw[answerred,line width=.6pt] (443.467,585.333)--(443.467,590.667);
\draw[answerred,line width=.6pt] (464.000,585.333)--(464.000,590.667);
\draw[answerred,line width=.6pt] (361.333,614.667)--(361.333,620.000);
\draw[answerred,line width=.6pt] (383.333,614.667)--(383.333,620.000);
\draw[answerred,line width=.6pt] (405.333,614.667)--(405.333,620.000);
\draw[answerred,line width=.6pt] (427.333,614.667)--(427.333,620.000);
\draw[answerred,line width=.6pt] (449.333,614.667)--(449.333,620.000);
\draw[answerred,line width=.6pt] (471.333,614.667)--(471.333,620.000);
\end{scope}
\end{scope}\pgfresetboundingbox\path[use as bounding box] (328.33,454.67) rectangle (558.33,519.0);
\end{tikzpicture}}\par
\Answer{\Formula{\begin{aligned}[t]&\phantom{={}}(56+64)\times4\\&=120\times4\\&=480\text{（千米）}\end{aligned}}\par 答：两地相距480千米。}

\MethodTitle{方法二：分别计算两部分，再求和}
\Answer{不先合并每份的数量，而是分别求出两部分的总数，再相加。}
\Answer{\Formula{\begin{aligned}[t]&\phantom{={}}56\times 4+64\times 4\\
&=224+256\\
&=480\text{（千米）}\end{aligned}}}
\Answer{答：两地相距480千米。}
\Problem{③甲、乙两个工程队一起铺路，甲队每天铺140米，乙队每天铺150米，两队一起干了5天，一共铺路多少米？}



\MethodTitle{方法一}

\par\noindent\makebox[\linewidth][c]{
\begin{tikzpicture}[x=1bp,y=-1bp,color=black,draw=black,text=black]
\begin{scope}\clip (355.67,561.33) rectangle (556.33,633.33);
% 高清原题重排：PDF 26，原书 21。不含解答。
\hdhead{21}{解决问题——信息窗3}
\hdsection{1}{145}{直接写得数。}
\q{96}{192}{193}{13}{\((15+25)\times2=\)}
\q{350}{192}{193}{13}{\((80-20)\div3=\)}
\q{597}{192}{193}{13}{\(5\times(12+18)=\)}
\q{96}{236}{193}{13}{\(36\div(4+5)=\)}
\q{350}{236}{193}{13}{\((25+15)\div8=\)}
\q{597}{236}{193}{13}{\(8\times(30-15)=\)}
\q{96}{280}{193}{13}{\(60\div(15-9)=\)}
\q{350}{280}{193}{13}{\((36+24)\div5=\)}
\q{597}{280}{193}{13}{\(30\times(100\div5)=\)}
\hdsection{2}{334}{解决问题。（先画图整理条件和问题，再解答。）}
\q{95}{378}{31}{13.5}{①}
\q{128}{379}{716}{13.5}{芳芳每分钟走80米，欣欣每分钟走75米，两人同时从家出发，5分}
\q{128}{421}{716}{13.5}{钟后同时到达少年宫。两家相距多少米？}
\q{462}{475}{105}{13.5}{芳芳家}
\q{615}{475}{105}{13.5}{少年宫}
\q{746}{475}{105}{13.5}{欣欣家}
\begin{scope}[scale=.6666667,draw=black,line width=.65pt]
\draw (495,509)--(779,509);
\fill (495,509) circle (1.9) (649,509) circle (1.9) (779,509) circle (1.9);
\end{scope}
\q{95}{597}{31}{13.5}{②}
\q{128}{598}{716}{13.5}{甲、乙两辆汽车同时从两地相对开出，甲车速度为56千米/时，乙车}
\q{128}{640}{716}{13.5}{速度为64千米/时，4小时后相遇。两地相距多少千米？}



\q{95}{972}{31}{13.5}{④}
\q{128}{973}{716}{13.5}{两艘船分别从A、B两港口相对开出，甲船的速度为32千米/时，乙}
\q{128}{1015}{716}{13.5}{船的速度为28千米/时，行驶5小时后还相距80千米。A、B两港口}
\q{128}{1057}{716}{13.5}{相距多少千米？}

\begin{scope}[color=answerred,draw=answerred,text=answerred]\ans{165.33}{137.33}{62.00}{9}{\linespread{1}\selectfont\setlength{\jot}{-2pt}\(\begin{aligned}[t]&=40\times2\\&=80\end{aligned}\)}
\ans{333.33}{137.33}{62.00}{9}{\linespread{1}\selectfont\setlength{\jot}{-2pt}\(\begin{aligned}[t]&=60\div3\\&=20\end{aligned}\)}
\ans{504.67}{137.33}{62.00}{9}{\linespread{1}\selectfont\setlength{\jot}{-2pt}\(\begin{aligned}[t]&=5\times30\\&=150\end{aligned}\)}
\ans{165.33}{167.33}{62.00}{9}{\linespread{1}\selectfont\setlength{\jot}{-2pt}\(\begin{aligned}[t]&=36\div9\\&=4\end{aligned}\)}
\ans{333.33}{167.33}{62.00}{9}{\linespread{1}\selectfont\setlength{\jot}{-2pt}\(\begin{aligned}[t]&=40\div8\\&=5\end{aligned}\)}
\ans{504.67}{167.33}{62.00}{9}{\linespread{1}\selectfont\setlength{\jot}{-2pt}\(\begin{aligned}[t]&=8\times15\\&=120\end{aligned}\)}
\ans{165.33}{196.00}{62.00}{9}{\linespread{1}\selectfont\setlength{\jot}{-2pt}\(\begin{aligned}[t]&=60\div6\\&=10\end{aligned}\)}
\ans{333.33}{196.00}{62.00}{9}{\linespread{1}\selectfont\setlength{\jot}{-2pt}\(\begin{aligned}[t]&=60\div5\\&=12\end{aligned}\)}
\ans{504.67}{196.00}{62.00}{9}{\linespread{1}\selectfont\setlength{\jot}{-2pt}\(\begin{aligned}[t]&=30\times20\\&=600\end{aligned}\)}
\ans{86.00}{308.67}{223.33}{11}{\linespread{1}\selectfont\setlength{\jot}{-2pt}\(\begin{aligned}[t]&\phantom{={}}(80+75)\times5\\&=155\times5\\&=775\text{（米）}\end{aligned}\)\\答：两家相距775米。}
\draw[answerred,line width=.9pt,->] (331.333,339.333)--(430.000,339.333);
\draw[answerred,line width=.9pt,->] (519.333,339.333)--(430.667,339.333);
\ans{350.67}{350.67}{80.67}{10}{400米}
\ans{454.00}{350.67}{73.33}{10}{375米}
\ans{330.00}{368.00}{212.00}{9.5}{芳芳：80米/分；欣欣：75米/分}
\ans{86.00}{451.33}{229.33}{11}{\linespread{1}\selectfont\setlength{\jot}{-2pt}\(\begin{aligned}[t]&\phantom{={}}(56+64)\times4\\&=120\times4\\&=480\text{（千米）}\end{aligned}\)\\答：两地相距480千米。}
\draw[answerred,line width=.8pt,->] (340.667,480.000)--(427.780,480.000);
\draw[answerred,line width=.8pt,->] (527.333,480.000)--(427.780,480.000);
\ans{333.33}{462.67}{37.33}{10}{甲}
\ans{520.00}{462.67}{33.33}{10}{乙}
\ans{414.00}{460.67}{52.67}{10}{相遇}
\ans{350.00}{489.33}{80.00}{9.5}{224千米}
\ans{448.00}{489.33}{80.00}{9.5}{256千米}
\ans{357.33}{504.00}{178.67}{9.5}{两车都行驶4小时}

\ans{360.67}{567.33}{170.67}{9.5}{甲：每天140米，铺5天}
\ans{360.67}{596.67}{170.67}{9.5}{乙：每天150米，铺5天}
\draw[answerred,line width=.8pt] (361.333,588.000)--(464.000,588.000);
\draw[answerred,line width=.8pt] (361.333,617.333)--(471.333,617.333);
\ans{473.33}{580.00}{74.00}{10}{700米}
\ans{480.00}{609.33}{71.33}{10}{750米}
\ans{86.00}{728.00}{245.33}{10.5}{\linespread{1}\selectfont\setlength{\jot}{-2pt}\(\begin{aligned}[t]&\phantom{={}}(32+28)\times5+80\\&=60\times5+80\\&=300+80\\&=380\text{（千米）}\end{aligned}\)\\答：两港口相距380千米。}
\draw[answerred,line width=.8pt] (343.333,751.333)--(536.667,751.333);
\draw[answerred,line width=.9pt,->] (343.333,751.333)--(424.737,751.333);
\draw[answerred,line width=.9pt,->] (536.667,751.333)--(465.439,751.333);
\ans{338.67}{733.33}{36.67}{10}{\(A\)}
\ans{529.33}{733.33}{30.00}{10}{\(B\)}
\ans{348.00}{760.00}{80.67}{9}{160千米}
\ans{425.33}{760.00}{48.67}{9}{80千米}
\ans{478.00}{760.00}{76.67}{9}{140千米}
\ans{388.00}{774.00}{168.67}{9}{两船各行驶5小时}
\draw[answerred,line width=.6pt] (361.333,585.333)--(361.333,590.667);
\draw[answerred,line width=.6pt] (381.867,585.333)--(381.867,590.667);
\draw[answerred,line width=.6pt] (402.400,585.333)--(402.400,590.667);
\draw[answerred,line width=.6pt] (422.933,585.333)--(422.933,590.667);
\draw[answerred,line width=.6pt] (443.467,585.333)--(443.467,590.667);
\draw[answerred,line width=.6pt] (464.000,585.333)--(464.000,590.667);
\draw[answerred,line width=.6pt] (361.333,614.667)--(361.333,620.000);
\draw[answerred,line width=.6pt] (383.333,614.667)--(383.333,620.000);
\draw[answerred,line width=.6pt] (405.333,614.667)--(405.333,620.000);
\draw[answerred,line width=.6pt] (427.333,614.667)--(427.333,620.000);
\draw[answerred,line width=.6pt] (449.333,614.667)--(449.333,620.000);
\draw[answerred,line width=.6pt] (471.333,614.667)--(471.333,620.000);
\end{scope}
\end{scope}\pgfresetboundingbox\path[use as bounding box] (355.67,561.33) rectangle (556.33,633.33);
\end{tikzpicture}}\par
\Answer{\Formula{\begin{aligned}[t]&\phantom{={}}(140+150)\times5\\&=290\times5\\&=1450\text{（米）}\end{aligned}}\par 答：一共铺路1450米。}

\MethodTitle{方法二：分别计算两部分，再求和}
\Answer{不先合并每份的数量，而是分别求出两部分的总数，再相加。}
\Answer{\Formula{\begin{aligned}[t]&\phantom{={}}140\times 5+150\times 5\\
&=700+750\\
&=1450\text{（米）}\end{aligned}}}
\Answer{答：一共铺路1450米。}
\Problem{④两艘船分别从A、B两港口相对开出，甲船的速度为32千米/时，乙船的速度为28千米/时，行驶5小时后还相距80千米。A、B两港口相距多少千米？}



\MethodTitle{方法一}

\par\noindent\makebox[\linewidth][c]{
\begin{tikzpicture}[x=1bp,y=-1bp,color=black,draw=black,text=black]
\begin{scope}\clip (333.67,727.33) rectangle (564.33,793.0);
% 高清原题重排：PDF 26，原书 21。不含解答。
\hdhead{21}{解决问题——信息窗3}
\hdsection{1}{145}{直接写得数。}
\q{96}{192}{193}{13}{\((15+25)\times2=\)}
\q{350}{192}{193}{13}{\((80-20)\div3=\)}
\q{597}{192}{193}{13}{\(5\times(12+18)=\)}
\q{96}{236}{193}{13}{\(36\div(4+5)=\)}
\q{350}{236}{193}{13}{\((25+15)\div8=\)}
\q{597}{236}{193}{13}{\(8\times(30-15)=\)}
\q{96}{280}{193}{13}{\(60\div(15-9)=\)}
\q{350}{280}{193}{13}{\((36+24)\div5=\)}
\q{597}{280}{193}{13}{\(30\times(100\div5)=\)}
\hdsection{2}{334}{解决问题。（先画图整理条件和问题，再解答。）}
\q{95}{378}{31}{13.5}{①}
\q{128}{379}{716}{13.5}{芳芳每分钟走80米，欣欣每分钟走75米，两人同时从家出发，5分}
\q{128}{421}{716}{13.5}{钟后同时到达少年宫。两家相距多少米？}
\q{462}{475}{105}{13.5}{芳芳家}
\q{615}{475}{105}{13.5}{少年宫}
\q{746}{475}{105}{13.5}{欣欣家}
\begin{scope}[scale=.6666667,draw=black,line width=.65pt]
\draw (495,509)--(779,509);
\fill (495,509) circle (1.9) (649,509) circle (1.9) (779,509) circle (1.9);
\end{scope}
\q{95}{597}{31}{13.5}{②}
\q{128}{598}{716}{13.5}{甲、乙两辆汽车同时从两地相对开出，甲车速度为56千米/时，乙车}
\q{128}{640}{716}{13.5}{速度为64千米/时，4小时后相遇。两地相距多少千米？}
\q{95}{783}{31}{13.5}{③}
\q{128}{784}{716}{13.5}{甲、乙两个工程队一起铺路，甲队每天铺140米，乙队每天铺150米，}
\q{128}{826}{716}{13.5}{两队一起干了5天，一共铺路多少米？}





\begin{scope}[color=answerred,draw=answerred,text=answerred]\ans{165.33}{137.33}{62.00}{9}{\linespread{1}\selectfont\setlength{\jot}{-2pt}\(\begin{aligned}[t]&=40\times2\\&=80\end{aligned}\)}
\ans{333.33}{137.33}{62.00}{9}{\linespread{1}\selectfont\setlength{\jot}{-2pt}\(\begin{aligned}[t]&=60\div3\\&=20\end{aligned}\)}
\ans{504.67}{137.33}{62.00}{9}{\linespread{1}\selectfont\setlength{\jot}{-2pt}\(\begin{aligned}[t]&=5\times30\\&=150\end{aligned}\)}
\ans{165.33}{167.33}{62.00}{9}{\linespread{1}\selectfont\setlength{\jot}{-2pt}\(\begin{aligned}[t]&=36\div9\\&=4\end{aligned}\)}
\ans{333.33}{167.33}{62.00}{9}{\linespread{1}\selectfont\setlength{\jot}{-2pt}\(\begin{aligned}[t]&=40\div8\\&=5\end{aligned}\)}
\ans{504.67}{167.33}{62.00}{9}{\linespread{1}\selectfont\setlength{\jot}{-2pt}\(\begin{aligned}[t]&=8\times15\\&=120\end{aligned}\)}
\ans{165.33}{196.00}{62.00}{9}{\linespread{1}\selectfont\setlength{\jot}{-2pt}\(\begin{aligned}[t]&=60\div6\\&=10\end{aligned}\)}
\ans{333.33}{196.00}{62.00}{9}{\linespread{1}\selectfont\setlength{\jot}{-2pt}\(\begin{aligned}[t]&=60\div5\\&=12\end{aligned}\)}
\ans{504.67}{196.00}{62.00}{9}{\linespread{1}\selectfont\setlength{\jot}{-2pt}\(\begin{aligned}[t]&=30\times20\\&=600\end{aligned}\)}
\ans{86.00}{308.67}{223.33}{11}{\linespread{1}\selectfont\setlength{\jot}{-2pt}\(\begin{aligned}[t]&\phantom{={}}(80+75)\times5\\&=155\times5\\&=775\text{（米）}\end{aligned}\)\\答：两家相距775米。}
\draw[answerred,line width=.9pt,->] (331.333,339.333)--(430.000,339.333);
\draw[answerred,line width=.9pt,->] (519.333,339.333)--(430.667,339.333);
\ans{350.67}{350.67}{80.67}{10}{400米}
\ans{454.00}{350.67}{73.33}{10}{375米}
\ans{330.00}{368.00}{212.00}{9.5}{芳芳：80米/分；欣欣：75米/分}
\ans{86.00}{451.33}{229.33}{11}{\linespread{1}\selectfont\setlength{\jot}{-2pt}\(\begin{aligned}[t]&\phantom{={}}(56+64)\times4\\&=120\times4\\&=480\text{（千米）}\end{aligned}\)\\答：两地相距480千米。}
\draw[answerred,line width=.8pt,->] (340.667,480.000)--(427.780,480.000);
\draw[answerred,line width=.8pt,->] (527.333,480.000)--(427.780,480.000);
\ans{333.33}{462.67}{37.33}{10}{甲}
\ans{520.00}{462.67}{33.33}{10}{乙}
\ans{414.00}{460.67}{52.67}{10}{相遇}
\ans{350.00}{489.33}{80.00}{9.5}{224千米}
\ans{448.00}{489.33}{80.00}{9.5}{256千米}
\ans{357.33}{504.00}{178.67}{9.5}{两车都行驶4小时}
\ans{86.00}{574.00}{244.00}{11}{\linespread{1}\selectfont\setlength{\jot}{-2pt}\(\begin{aligned}[t]&\phantom{={}}(140+150)\times5\\&=290\times5\\&=1450\text{（米）}\end{aligned}\)\\答：一共铺路1450米。}
\ans{360.67}{567.33}{170.67}{9.5}{甲：每天140米，铺5天}
\ans{360.67}{596.67}{170.67}{9.5}{乙：每天150米，铺5天}
\draw[answerred,line width=.8pt] (361.333,588.000)--(464.000,588.000);
\draw[answerred,line width=.8pt] (361.333,617.333)--(471.333,617.333);
\ans{473.33}{580.00}{74.00}{10}{700米}
\ans{480.00}{609.33}{71.33}{10}{750米}

\draw[answerred,line width=.8pt] (343.333,751.333)--(536.667,751.333);
\draw[answerred,line width=.9pt,->] (343.333,751.333)--(424.737,751.333);
\draw[answerred,line width=.9pt,->] (536.667,751.333)--(465.439,751.333);
\ans{338.67}{733.33}{36.67}{10}{\(A\)}
\ans{529.33}{733.33}{30.00}{10}{\(B\)}
\ans{348.00}{760.00}{80.67}{9}{160千米}
\ans{425.33}{760.00}{48.67}{9}{80千米}
\ans{478.00}{760.00}{76.67}{9}{140千米}
\ans{388.00}{774.00}{168.67}{9}{两船各行驶5小时}
\draw[answerred,line width=.6pt] (361.333,585.333)--(361.333,590.667);
\draw[answerred,line width=.6pt] (381.867,585.333)--(381.867,590.667);
\draw[answerred,line width=.6pt] (402.400,585.333)--(402.400,590.667);
\draw[answerred,line width=.6pt] (422.933,585.333)--(422.933,590.667);
\draw[answerred,line width=.6pt] (443.467,585.333)--(443.467,590.667);
\draw[answerred,line width=.6pt] (464.000,585.333)--(464.000,590.667);
\draw[answerred,line width=.6pt] (361.333,614.667)--(361.333,620.000);
\draw[answerred,line width=.6pt] (383.333,614.667)--(383.333,620.000);
\draw[answerred,line width=.6pt] (405.333,614.667)--(405.333,620.000);
\draw[answerred,line width=.6pt] (427.333,614.667)--(427.333,620.000);
\draw[answerred,line width=.6pt] (449.333,614.667)--(449.333,620.000);
\draw[answerred,line width=.6pt] (471.333,614.667)--(471.333,620.000);
\end{scope}
\end{scope}\pgfresetboundingbox\path[use as bounding box] (333.67,727.33) rectangle (564.33,793.0);
\end{tikzpicture}}\par
\Answer{\Formula{\begin{aligned}[t]&\phantom{={}}(32+28)\times5+80\\&=60\times5+80\\&=300+80\\&=380\text{（千米）}\end{aligned}}\par 答：两港口相距380千米。}

\MethodTitle{方法二：分别求两船已行路程，再加剩余距离}
\Answer{两港距离包括甲船行程、乙船行程和中间剩余的80千米。}
\Answer{\Formula{\begin{aligned}[t]&\phantom{={}}32\times 5+28\times 5+80\\
&=160+140+80\\
&=300+80\\
&=380\text{（千米）}\end{aligned}}}
\Answer{答：两港口相距380千米。}
